\documentclass[11pt]{article}
\usepackage[a4paper,margin=1in]{geometry}
\usepackage{amsmath,amssymb,amsfonts}
\usepackage{hyperref}
\usepackage{graphicx}
\usepackage{booktabs}
\usepackage{microtype}
\usepackage{siunitx}
\usepackage[numbers,sort&compress]{natbib}
\usepackage{algorithm}
\usepackage{algpseudocode}
\usepackage{listings}
\usepackage{xcolor}

\hypersetup{colorlinks=true, linkcolor=blue, citecolor=blue, urlcolor=blue}

\title{M--QNN v2: A NISQ-Practical Quantum Classifier with Classical-Shadow Screening and Adaptive Shot Allocation}
\author{Ryan O. Van Gelder \\ \\ Citizen Gardens}
\date{October 23, 2025}

\lstdefinestyle{code}{
  basicstyle=\ttfamily\small,
  breaklines=true,
  frame=single,
  numbers=left,
  numbersep=5pt,
  xleftmargin=2em,
  framexleftmargin=1.5em,
  language=Python
}

\begin{document}
\maketitle

\begin{abstract}
We formalize a NISQ-practical quantum classification pipeline that combines a variational encoder, soft $k$-nearest neighbors (k-NN) using fidelity distances, classical-shadow screening to reduce candidate neighbors, and adaptive shot allocation with certified early-exit. Optional zero-noise extrapolation (ZNE) improves measurement robustness. We provide algorithms, complexity accounting, and reference implementations. We make no quantum-advantage claims.
\end{abstract}

\section{Scope and Notation}
We target pure-state encodings on $n$ qubits. Let $\theta \in \mathbb{R}^p$ denote encoder parameters and $x \in \mathbb{R}^d$ be features. The encoder prepares
\[
|\psi_\theta(x)\rangle = U_\theta(\Phi(x))|0\rangle^{\otimes n},
\]
with optional data re-uploading \citep{PerezSalinas2020}. Labels are in $\mathcal{Y}=\{1,\dots,C\}$ over a dataset $\mathcal{D}=\{(x_i,y_i)\}_{i=1}^N$. For pure states, fidelity is $F(|\psi\rangle,|\phi\rangle)=|\langle\psi|\phi\rangle|^2 \in [0,1]$, and we use distance $d_F=1-F$.

\section{Model}
\subsection{Variational encoder}
We use hardware-efficient ans\"atze with shallow depth and native entanglers \citep{Kandala2017}. Gradients are estimated via parameter-shift rules \citep{Schuld2019,Wierichs2022}.

\subsection{Soft k-NN with classical-shadow screening}
We precompute $m$ single-shot global Clifford measurements (``classical shadows'') for each reference state and for the query state, using shared randomness. Classical shadows allow accurate post-hoc prediction of many state properties, including fidelities, with logarithmic sample scaling \citep{Huang2020,ClassicalShadowsNoise}. We compute a proxy similarity $S_{\mathrm{proxy}}(x,x_j)$ from inner products of shadow features and keep the top $M \ll |\mathcal{R}|$ candidates. Exact fidelities are then measured only on the candidate set via the swap test \citep{Buhrman2001}.

Soft weights for class probabilities are
\begin{equation}
w_j(x)=\frac{\exp\{-\tau\,(1-F(x,x_j))\}}{\sum_{\ell\in\mathcal{C}(x)} \exp\{-\tau\,(1-F(x,x_\ell))\}},\quad
\hat p_c(x)=\sum_{j\in\mathcal{C}(x)} w_j(x)\,\mathbf{1}[y_j=c],
\end{equation}
with temperature $\tau>0$ and candidate set $\mathcal{C}(x)$. Training minimizes cross-entropy over minibatches and backpropagates through the encoder via parameter shift.

\subsection{Adaptive shots with certified early-exit}
Let $\hat F_j$ be the current estimate after $S_j$ shots for neighbor $j$. Using Hoeffding's inequality \citep{Hoeffding1963}, define a confidence radius $r_j=\sqrt{\ln(2/\delta)/(2S_j)}$ for failure probability $\delta$. Stop when a winner $j^\star$ satisfies
\begin{equation}
\hat F_{j^\star}-r_{j^\star} > \max_{j\neq j^\star} (\hat F_j + r_j).
\end{equation}
Allocate the next shot to the index with the largest uncertainty or smallest margin to improve stopping efficiency.

\section{Fidelity estimation via the swap test}
Preparing $|\psi_\theta(x)\rangle$ and $|\psi_\theta(x_j)\rangle$, the swap test returns
\begin{equation}
\Pr[\text{ancilla}=0] = \tfrac{1}{2}\left(1 + F(x,x_j)\right),\qquad \hat F=2\hat p_0-1.
\end{equation}
The estimator variance obeys $\mathrm{Var}(\hat F)\le 1/S$ under i.i.d.\ shots.

\section{Zero-Noise Extrapolation (optional)}
We scale noise by circuit folding and extrapolate to zero noise \citep{Temme2017,He2020,GiurgicaTiron2020}. For scale factors $\lambda\in\{1,3,5\}$, fit a linear model $\hat F(\lambda)\approx a_0 + a_1(\lambda-1)$ and take $\hat F_{\mathrm{zne}}=a_0$. Policy: reject extrapolation if $R^2<0.9$ or $|\hat F_{\mathrm{zne}}-\hat F_{\lambda=1}|>0.2$; cap folded depth within coherence.

\section{Training algorithm}
\textbf{Inputs:} $\mathcal{D}$, reference set $\mathcal{R}$, $\tau$, $\eta$, shadows $m$, candidates $M$, confidence $\delta$, refresh period $K$.

\begin{algorithm}[H]
\caption{Enhanced training with screening and stochastic parameter-shift}
\begin{algorithmic}[1]
\State Initialize $\theta$
\State Precompute $m$ shadows for all references in $\mathcal{R}$; refresh every $K$ epochs
\For{epochs}
  \For{minibatch $B\subset \mathcal{D}$}
    \For{sample $x\in B$}
      \State Prepare $|\psi_\theta(x)\rangle$ and collect $m$ shadows
      \State Select candidates $\mathcal{C}(x)$ of size $M$ using proxy distances
      \For{$x_j\in\mathcal{C}(x)$}
        \State Estimate $F(x,x_j)$ with adaptive shots and early-exit
      \EndFor
      \State Compute soft weights $w_j(x)$ and probabilities $\hat p(x)$
    \EndFor
    \State Update $\theta$ by stochastic parameter-shift on a random subset $P_t\subset\{1,\dots,p\}$
  \EndFor
\EndFor
\end{algorithmic}
\end{algorithm}

\section{Resource and complexity control}
Swap-test work scales as $\approx |B|\,M\,\bar S$ per epoch after screening; encoder evaluations scale as $\approx 2|P_t|\,|B|$ per epoch for gradients. Shadow overhead is $\approx m(|B|+|\mathcal{R}|)$ per refresh. Recommended settings: $m\in[64,256]$, $M\in[16,64]$, $\delta\in[0.01,0.05]$, $|P_t|/|P|\in[0.1,0.3]$.

\section{Evaluation protocol}
Datasets: synthetic $2$--$4$\,qubit oracles and a small real dataset mapped to $\mathbb{R}^d$. Baselines: classical k-NN and, for synthetic data, exact fidelities computed classically. Metrics: accuracy, macro-F1, calibration (ECE) with 95\% CIs via bootstrap and shot resampling. Ablations: screening recall, early-exit shot savings, stochastic shift variance, coreset medoids \citep{KaufmanRousseeuw1990}.

\section{Practical constraints and risks}
Pure-state assumption for the swap test; mixed states are out of scope. Depth bounded by coherence and native two-qubit error. Barren plateaus are a risk for deep encoders \citep{McClean2018}; prefer shallow, local entanglers and monitor gradient norms. Data movement is non-trivial; exploit batching and reuse if the platform allows.

\section{Code snippets}
\subsection*{Qiskit: swap-test fidelity}
\begin{lstlisting}[style=code,caption={Swap-test fidelity in Qiskit},label={lst:swaptest}]
from qiskit import QuantumCircuit, QuantumRegister, ClassicalRegister, Aer, transpile
from qiskit.quantum_info import Statevector

def swap_test_circuit(psi: QuantumCircuit, phi: QuantumCircuit):
    n = psi.num_qubits
    anc = QuantumRegister(1, "anc")
    a = QuantumRegister(n, "a")
    b = QuantumRegister(n, "b")
    c = ClassicalRegister(1, "c")
    qc = QuantumCircuit(anc, a, b, c)

    qc.h(anc)                     # prepare |+>
    qc.compose(psi, qubits=a, inplace=True)
    qc.compose(phi, qubits=b, inplace=True)

    # controlled-SWAP between registers a and b
    for i in range(n):
        qc.cswap(anc[0], a[i], b[i])

    qc.h(anc)
    qc.measure(anc, c)
    return qc

def estimate_fidelity(psi, phi, shots=1000, backend=None):
    backend = backend or Aer.get_backend("aer_simulator")
    qc = swap_test_circuit(psi, phi)
    t_qc = transpile(qc, backend)
    result = backend.run(t_qc, shots=shots).result()
    counts = result.get_counts()
    p0 = counts.get('0', 0) / shots
    return 2*p0 - 1  # unbiased estimator for F
\end{lstlisting}

\subsection*{Adaptive early-exit}
\begin{lstlisting}[style=code,caption={Hoeffding-interval early-exit},label={lst:earlyexit}]
import math

def should_stop(winner_F, winner_S, others, delta=0.01):
    # others: list of tuples (F_j, S_j)
    r_star = math.sqrt(math.log(2/delta) / (2*max(1, winner_S)))
    lhs = winner_F - r_star
    rhs = max(F + math.sqrt(math.log(2/delta) / (2*max(1, S))) for F, S in others)
    return lhs > rhs
\end{lstlisting}

\subsection*{PennyLane: parameter-shift gradient skeleton}
\begin{lstlisting}[style=code,caption={Parameter-shift training skeleton},label={lst:psr}]
import pennylane as qml
import numpy as np

n_qubits = 4
dev = qml.device("default.qubit", wires=n_qubits)

def encoder(x, params):
    # simple hardware-efficient layer
    for i in range(n_qubits):
        qml.RZ(params[0, i]*x[0] + params[1, i], wires=i)
        qml.RX(params[2, i]*x[0] + params[3, i], wires=i)
    for i in range(n_qubits-1):
        qml.CNOT(wires=[i, i+1])

@qml.qnode(dev, interface="autograd")
def state_overlap_cost(x, xj, params):
    # prepare |psi_theta(x)> and |psi_theta(xj)> on two registers if needed,
    # or encode a proxy observable for training the encoder
    encoder(x, params)
    # return expectation of a simple observable as a surrogate loss
    return qml.expval(qml.PauliZ(0))

# autograd/parameter-shift is used under the hood
params = 0.01*np.random.randn(4, n_qubits)
x = np.array([0.5])
xj = np.array([0.2])
grad = qml.grad(lambda p: state_overlap_cost(x, xj, p))(params)
\end{lstlisting}

\section{Reproducibility}
Report full $U_\theta$ and $\Phi(x)$, depth, gate counts, screening $(m,M)$, $\delta$, $|P_t|/|P|$, coreset sizes, ZNE policy, raw and ZNE-adjusted fidelities with CIs, early-exit statistics, seeds, and code.

\bibliographystyle{unsrtnat}
\bibliography{references}
\end{document}
